% retained preamble for sources/alco-235-vertex-models-for-canonical-grothendieck-polynomials-2023
% Assembled by the source-intake line 2026-10-15, per the consumer's six-round error chain spec:
%   mathtools + \DeclarePairedDelimiter, \address, \datereceived/\dateaccepted/\daterevised as no-ops,
%   the abstract environment family, and the ALCO theorem family; author macros kept verbatim.
% Comments are ASCII-only and no orphan metadata arguments remain (delivery rule).
\PassOptionsToPackage{dvipsnames,svgnames,table}{xcolor}
\documentclass[11pt]{article}
\usepackage[a4paper,margin=22mm]{geometry}
\usepackage[no-math]{fontspec}
\setmainfont{DejaVu Serif}
\usepackage{amsmath,amssymb,amsthm,mathtools,graphicx,xcolor,hyperref,xurl,xspace,ifthen,enumerate}
\usepackage[shortlabels]{enumitem}
\setlength{\emergencystretch}{4em}
% ---- publisher-only commands as no-ops (cedram class) ----
\providecommand{\address}[1]{}
\providecommand{\email}[1]{}
\providecommand{\urladdr}[1]{}
\providecommand{\datereceived}[1]{}
\providecommand{\dateaccepted}[1]{}
\providecommand{\daterevised}[1]{}
\providecommand{\datepublished}[1]{}
\providecommand{\subjclass}[1]{}
\providecommand{\keywords}[1]{}
% ---- abstract environment family ----
\providecommand{\abstractname}{Abstract}
\ifdefined\abstract\else\newenvironment{abstract}{\begin{quote}\small\noindent\textbf{\abstractname.}\ }{}\fi
% ---- ALCO theorem family (shared counter, publisher numbering) ----
\newtheorem{theo}{Theorem}[section]
\newtheorem{lemm}[theo]{Lemma}
\newtheorem{coro}[theo]{Corollary}
\newtheorem{prop}[theo]{Proposition}
\newtheorem{defi}[theo]{Definition}
\newtheorem{rema}[theo]{Remark}
\newtheorem{exam}[theo]{Example}
\newtheorem{ques}[theo]{Question}
\newtheorem{conj}[theo]{Conjecture}
% ---- author macros (verbatim) ----
%\usepackage{fourier}
\usepackage[T1]{fontenc}
%\usepackage[a4paper,left=2.25cm,right=2.25cm,top=3cm,bottom=3cm,headsep=1cm]{geometry}
%\setlength{\parskip}{1pt}
\usepackage{hyperref}
\usepackage{cleveref}
%\usepackage{mathtools} %for automatic size of the brackets
%tikz stuff
\usepackage{tikz}\usetikzlibrary{graphs,quotes,fit,positioning,matrix,calc,decorations.markings,angles,decorations.pathmorphing,decorations.pathreplacing}%,matrix,fit,shapes.geometric}
% useful macros
\tikzset{math mode/.style = {execute at begin node=$, execute at end node=$}}%makes all nodes math mode
\tikzset{arrow/.style={postaction={decorate,thick,decoration={markings,mark = at position #1 with {\arrow{>}}}}},arrow/.default=0.5}
\tikzset{invarrow/.style={postaction={decorate,thick,decoration={markings,mark = at position #1 with {\arrow{<}}}}},invarrow/.default=0.5}
\newcommand\bbull[3]{\filldraw[fill=black!35!blue, draw=black] (#1,#2) circle (#3cm);}
\newcommand\ebull[3]{\filldraw[fill=white,draw=black] (#1,#2) circle (#3cm);}
\tikzset{fermionic/.style={thin}}
\tikzset{bosonic/.style={ultra thick}}
\tikzset{gcol/.style={Bo}}
\tikzset{gccol/.style={red}}
\tikzset{Gcol/.style={blue}}
% 
\DeclarePairedDelimiter\autobracket{(}{)}
\newcommand{\br}[1]{\autobracket*{#1}}
%
\definecolor{Bo}{rgb}{1.0, 0.49, 0.0}
\definecolor{Bs}{rgb}{1.0, 0.44, 0.37}
\definecolor{aj}{rgb}{0.8, 0.8, 1.0}
%
\newtheorem{case}{Case}
\newcommand{\sslash}{\mathbin{/\mkern-6mu/}}
\newcommand{\bra}[1]{\left\langle #1\right|}
\newcommand{\ket}[1]{\left|#1\right\rangle}
\newcommand{\braket}[2]{\left< #1 \vphantom{#2}\right|
\! \left. #2 \vphantom{#1} \right>}%
%\usepackage[toc,page]{appendix}
%\usepackage[inline]{showlabels}
\setlength{\multlinegap}{0pt}
\allowdisplaybreaks[1]
%
\renewcommand\ss{\scriptstyle}
\newcommand\sss{\scriptscriptstyle}
\newcommand\ds{\displaystyle}
\def\phid{\phi^\dagger}
\def\psid{\psi^\dagger}
\def\fs{\footnotesize}
\newcommand\g{{g^{(\alpha,\beta)}_{\lambda}}}
\newcommand\Gg{{G^{(\alpha,\beta)}_{\lambda}}}
\newcommand\lb{{\left(}}
\newcommand\rb{{\right)}}
\newcommand\ZZ{{\mathbb Z}}
\newcommand\QQ{{\mathbb Q}}
\newcommand\CC{{\mathbb C}}
\newcommand\PP{{\mathbb P}}
\newcommand\R{{\mathcal {R}}}
\newcommand\LL{{\mathcal {L}}}
\newcommand\jl{{\mathfrak l}}
\newcommand\jt{{\mathfrak t}}


\long\def\junk#1{}







% The abstract text of the source is dropped here: it belongs to the document body, not to the preamble, and the assembler supplies the front matter itself.


% Added for this card: the retained statement uses the thm environment, which the delivered preamble does not declare, so it is aliased to the shared theorem counter.
\newtheorem{thm}[theo]{Theorem}
\crefname{originalequation}{original equation}{original equations}
\Crefname{originalequation}{Original equation}{Original equations}
\makeatletter
\newlabel{branchG}{{2.13}{source}{}{sourcepointer-1}{}}
\newlabel{branchG@cref}{{[originalequation][0][]2.13}{[1][0][]source}{}{}{}}
\makeatother
\title{The canonical Grothendieck polynomials are given by the row multiplicity vertex model}
\author{Ajeeth Gunna, Paul Zinn-Justin}
\date{doi:10.5802/alco.235; CC BY 4.0}
\begin{document}
\maketitle
\setcounter{section}{2}
\setcounter{theo}{1}
\noindent\textit{Editorial note.} The counted unit is the article's Theorem thm:rowG at source lines 733-758 together with its proof at lines 760-790, an adjacent pairing read from the source. The statement is printed 2.2: the number is taken from the census adjudication and the counters are seeded accordingly. The source is the publisher version of the article in Algebraic Combinatorics, DOI 10.5802/alco.235, released under CC BY 4.0.\par\medskip
\noindent\textit{Editorial note.} The retained spans contain no citation and no reference to another statement: the two equation labels they use are defined inside the statement itself, so no pointer and no equation stub is needed.\par\medskip
\noindent\textit{Editorial note.} The source ships no retained preamble, and the environments of the retained span come from the publisher class, which the assembler does not carry; the card therefore slices the preamble from the pinned source and adds the theorem declarations, recording that addition here. The unit does ship the retained bibliography prepared by the intake line, which this face does not need.\par\medskip
\noindent\textit{Editorial note.} This is the first face of the unit; the unit had been frozen after it pending the Lead ruling and was released to this lane, so further faces of the same unit may follow.\par\medskip
\noindent\textit{Editorial note.} Revision r002 drops the input lines from the sliced preamble: the assembler requires the card to carry the text of every passage it uses rather than pulling files at build time. Dropped line kinds: ['\textbackslash{}the inclusion of the tableau macros']\par\medskip
\noindent\textit{Editorial note.} Revision r003 lists the mathtools package: the sliced preamble uses the paired-delimiter declaration that the publisher class loads through the input file this card had to drop, so the package is named explicitly instead.\par\medskip
\noindent\textit{Editorial note.} Revision r004 drops the publisher metadata commands from the sliced preamble (\textbackslash{}title\{Vertex models for Can, \textbackslash{}author\{\textbackslash{}firstname\{Ajeeth\} \textbackslash{}, \textbackslash{}address\{School of Mathemati, \textbackslash{}email\{agunna@student.unimel, \textbackslash{}author\{\textbackslash{}firstname\{Paul\} \textbackslash{}la, \textbackslash{}address\{School of Mathemati); the assembler supplies the title and the author block, and the class that defines those commands is not available to the build.\par\medskip
\noindent\textit{Editorial note.} Revision r006 defines the publisher date commands as no-ops, since the class that defines them is not available to the build. Commands covered: \textbackslash{}dateaccepted, \textbackslash{}datereceived, \textbackslash{}daterevised.\par\medskip
\noindent\textit{Editorial note.} Revision r007 uses the retained preamble delivered by the intake line for this unit, which settles all six items the earlier revisions exposed one by one: the packages including mathtools, the no-ops for the publisher metadata and date commands, the abstract environment, and the theorem family under the shared counter. It is used as shipped, with one alias added for the theorem environment of the retained statement.\par\medskip
\noindent\textit{Editorial note.} Revision r008 drops the file-inclusion line that the delivered preamble still carried: the assembler requires every passage to be flattened into the card, and that included file is not shipped with the source, so the line cannot be honoured. The task notes were also reworded so that no note quoted the inclusion command itself, which the same check flags.\par\medskip
\noindent\textit{Editorial note.} Revision r009 removes the abstract text the delivered preamble still carried at its end (4 lines): that text belongs to the body of the article, while the assembler supplies the front matter of the excerpt, and a document environment opening inside the preamble stops the build.\par\medskip
\noindent\textit{Editorial note.} Revision r010 removes the publisher metadata commands and their arguments that the delivered preamble still carried at the point of use (subject classification, keywords and the date commands): the assembler supplies the front matter and the class defining those commands is not available to the build. Lines touched: 1.\par\medskip
\noindent\textit{Editorial note.} Revision r011 removes the last orphan of the publisher front matter: one line of the article title had stayed behind after the title command itself was dropped, and a bare text line inside the preamble stops the build. Dropped: Exactly solvable lattice models\}.\par\medskip
\noindent\textit{Editorial note.} Revision r012 seeds the shared theorem counter theo instead of thm: the delivered preamble aliases the theorem environment onto that counter, so the name thm is not a counter at all and the earlier seeding stopped the build.\par\medskip
\noindent\textit{Editorial note.} Revision r013 declares the one reference the retained proof makes to a figure outside the excerpt, which the previous build reported as undefined.\par\medskip
\noindent\textit{Editorial note.} Revision r014 corrects that declaration: the reference is to an equation of the article, not to a figure, and its number is computed from the source rather than assumed.\par\medskip
\noindent\textit{Allowed standard prerequisites (editorial).}\begin{itemize}
\item Vertex model for partitions recorded with row multiplicities. The article encodes a partition by the multiplicities of its rows and builds a vertex model on a grid whose vertices carry labels, so that the weight of a configuration reproduces the canonical Grothendieck polynomial; the matrix element of the transfer matrix counts the configurations compatible with the boundary data. The author's own definitions are the ones used by the counted statement and proof.. The counted theorem identifies the polynomial with the row multiplicity vertex model, which the article then uses to read off combinatorial properties of the coefficients.
\end{itemize}\medskip
\ifdefined\hypertarget\hypertarget{sourcepointer-1}{}\fi
\noindent\textit{Source reference.} Original equation 2.13: sources/alco-235-vertex-models-for-canonical-grothendieck-polynomials-2023/pinned-original-source.tex, source lines 660--672. Navigation pointer invoked by the retained proof when it compares the branching formula for the polynomial family: that formula is an equation of the article that lies outside the excerpt. The reference is declared as a pointer carrying the number computed from the source by counting the rows of its numbered display environments section by section, since this article resets its equation counter at every section: it is equation 2.13.\par\medskip
\begin{thm}
\label{thm:rowG}
The canonical Grothendieck polynomials $\Gg(x)$  are given by
%
\begin{align}
\label{J}
\Gg(x_1,\dots,x_n)
&=
\bra{0}
T(x_1)
\dots
T(x_n)
\ket{\lambda}
\\
\label{dJ}
\Gg(x_1,\dots,x_n)
&=
\bra{\lambda}
T^{*}(x_n)
\dots
T^{*}(x_1)
\ket{0}
\end{align}
%
where $\ket{\lambda} = \bigotimes_{i=1}^{\infty} \ket{m_i(\lambda)}$, and similarly for the dual state $\bra{\lambda}$.
\end{thm}

\begin{proof}
We shall prove \eqref{J}, and \eqref{dJ} follows immediately as a consequence of the way we defined the $L^{*}$ matrix. Fix $\lambda$, then we can just consider the finite transfer matrix of size $\lambda_1$. Inserting a complete set of states before the final transfer matrix, we have the following branching formula

\begin{align*}
    \Gg(x_1,\dots,x_n,x)
&=
\sum_{\mu}
\bra{0}
T(x_1)
\dots
T(x_n)
\ket{\mu}
\bra{\mu}
T(x)
\ket{\lambda}.
\end{align*}

On comparing the branching formula for $\Gg$ (\cref{branchG}), it is enough to show $G^{(\alpha,\beta)}_{\lambda\sslash\mu}(x)=\bra{\mu}T(x)\ket{\lambda}$.
Recall that for a horizontal strip $\lambda/\mu$, we have
\begin{align*}
G^{(\alpha,\beta)}_{\lambda\sslash\mu}(x)= 
\left(
\frac{x}{1-\alpha x}
\right)^{|\lambda/\mu|}
\left(
\frac{1+\beta x}{1-\alpha x}
\right)^{r(\mu/\Bar{\lambda})}.
\end{align*}

Based on the vertices that we used to define $T$, one easily observes that $\bra{\mu}T(x) \ket{\lambda}\neq 0$ if and only if $\lambda/\mu$ is a horizontal strip. The label $1$ on the left edge at site $i$ amounts to adding a box in the $i^{th}$ column from the left. For every such vertex, we get a factor of $\frac{x}{1-\alpha x}$. From our previous analysis, we see that $r(\mu/\Bar{\lambda})$ is exactly the number of vertices with the label $0$ on the left edge and a non zero label on the bottom edge, where each such vertex has a weight of $\frac{1+\beta x}{1-\alpha x}$.
\end{proof}

\end{document}
